\label{sec:large-scale-experiment}
The previous section evaluates \toolname\ under established benchmark protocols
used by prior invariant and specification-generation tools. Many of those
programs include explicit \texttt{assert} statements, which are useful for prior tools to generate specifications strong enough to discharge those assertions, i.e., verification goals, but they do not fully exercise the goal-independent setting \toolname\ is
designed to support. To complement that evaluation, we run a separate
study on goal-independent programs, where all \texttt{assert} statements are
removed and the generator must decide what functional behavior to specify.
For this comparison we select the state-of-the-art specification-generation
tool in each of the two language ecosystems we span: \autospec~\cite{autospec}
for C/ACSL and \specgen~\cite{SpecGen} for Java/JML.

\subsection{Experimental Setup}
We construct \textbf{\toolname-400}, a 400-program C benchmark built from three specification-generation suites:
\textbf{FormalBench}~\cite{le2025can} (250 programs) and \textbf{SpecGenBench} (59), translated to C from Java, and the native-C \textbf{SV-COMP} (91); programs are stripped of \texttt{assert} statements so
the generator must decide what to specify. The \textbf{FormalBench} base dataset contains
Java programs verified in OpenJML with sequential, branching, and loop-based
structures. The translation preserves the source-level control flow and data
dependencies used by the benchmark programs. Verification
follows the Frama-C/WP protocol of Section~\ref{sec:experiments} (Z3, Typed
memory model).

All LLM-based methods are evaluated on \texttt{gpt-4o}, \texttt{gpt-5-mini},
\texttt{gpt-5.4-mini}, and \texttt{gpt-5} at sampling temperature 0.7 under
Pass@1. We deliberately draw all four from a single model family: models from
different families differ too much in training data and capability for a clean
comparison, whereas staying within one family lets performance differences
reflect reasoning capability rather than cross-family data bias. This set also
differs from the models in Section~\ref{sec:experiments}, which were fixed by the
baselines compared there.

We compare four C/ACSL workflows: full \toolname\ (up to 10 refinement
iterations), \toolname\ \emph{w/o refine}, \toolname\ \emph{w/o refine w/o
symbolic execution (Pure LLM)}, and \autospec\ in its default one-shot
configuration; plus \specgen, which is run on the original Java
versions in its native Java/JML pipeline. The two ablations, \emph{Pure LLM} and
\emph{w/o refine}, isolate the contributions of symbolic-execution guidance and
verifier-guided refinement by comparison with the full pipeline, while \autospec
serves as the verifier-level baseline that shares our C/ACSL setting. \specgen\ is
the closest goal-independent system but runs in the Java/JML ecosystem, so we
include it only as a cross-ecosystem reference rather than a direct
verifier-level baseline. Each run is capped at a 1800-second outer timeout.

This study, run on all 400 \textbf{\toolname-400} programs without the per-program compute cap of Section~\ref{sec:experiments}, probes the goal-independent setting along four
orthogonal axes: whether the pipeline produces \emph{valid} specifications, whether
those specifications are \emph{strong}, at what \emph{cost}, and how it
\emph{fails}. Throughout, we say a program is \emph{valid} when the
generated specification is accepted in full by the verifier, i.e., \emph{all}
verification conditions it induces are discharged within the configured budget; a
specification with any unproved obligation is not counted, so validity here measures
whole-specification verification rather than partial proof.
Extending the four research questions of Section~\ref{sec:experiments}, this
study is organized around four further questions, one per axis:
\begin{itemize}
    \item \textbf{RQ5:} How often does \toolname\ produce a
    \emph{Valid} specification?
    \item \textbf{RQ6:} How strong are \toolname's specifications
    versus the baselines?
    \item \textbf{RQ7:} What does \toolname\ cost in time, tokens, and
    API calls?
    \item \textbf{RQ8:} What are \toolname's dominant failure modes?
\end{itemize}

\subsection{RQ5: How often does \toolname\ produce a \emph{Valid} specification?}
\label{subsec:rq5-valid}

% Auto-generated by src/scripts/aggregate_all_methods_all_models.py
\begin{table}[t]
\centering
\small
\caption{Native-verifier validity on \toolname-400.}
\label{tab:large_experiment_summary}
\setlength{\tabcolsep}{3.5pt}
\begin{tabular}{llrr}
\toprule
Method & Model & Valid & Valid (\%) \\
\midrule
\multirow{4}{*}{\autospec\ (C/ACSL)}
 & gpt-4o & 115 & 28.8 \\
 & gpt-5-mini & 104 & 26.0 \\
 & gpt-5.4-mini & 255 & 63.8 \\
 & gpt-5 & 288 & 72.0 \\
\midrule
\multirow{4}{*}{\specgen\ (Java/JML)}
 & gpt-4o & 144 & 36.0 \\
 & gpt-5-mini & 161 & 40.2 \\
 & gpt-5.4-mini & 198 & 49.5 \\
 & gpt-5 & 251 & 62.7 \\
\midrule
\multirow{4}{*}{\shortstack[l]{\toolname\\w/o refine\\w/o symbolic execution (Pure LLM)}}
 & gpt-4o & 52 & 13.0 \\
 & gpt-5-mini & 76 & 19.0 \\
 & gpt-5.4-mini & 106 & 26.5 \\
 & gpt-5 & 55 & 13.8 \\
\midrule
\multirow{4}{*}{\toolname\ w/o refine}
 & gpt-4o & 126 & 31.5 \\
 & gpt-5-mini & 107 & 26.8 \\
 & gpt-5.4-mini & 132 & 33.0 \\
 & gpt-5 & 144 & 36.0 \\
\midrule
\multirow{4}{*}{\textbf{\toolname}}
 & gpt-4o & 272 & 68.0 \\
 & gpt-5-mini & 333 & 83.3 \\
 & gpt-5.4-mini & 373 & 93.3 \\
 & gpt-5 & 372 & 93.0 \\
\bottomrule
\end{tabular}
\end{table}


Table~\ref{tab:large_experiment_summary} reports the per-model results. Within the C/ACSL workflows, the full \toolname\ pipeline achieves the highest \emph{Valid} rate on every model: \(68.0\%\) on gpt-4o, \(83.3\%\) on gpt-5-mini, \(93.3\%\) on gpt-5.4-mini, and \(93.0\%\) on gpt-5. These correspond to absolute gains of \(+55.0\), \(+64.3\), \(+66.8\), and \(+79.2\)\,pp over Pure LLM, and \(+39.2\), \(+57.3\), \(+29.5\), and \(+21.0\)\,pp over \autospec.
The ablations show two different sources of improvement. Refinement is the dominant source of validity gain, but it is also the most expensive component: enabling refinement lifts \emph{Valid} by \(+36.5\), \(+56.5\), \(+60.3\), and \(+57.0\)\,pp across the four models, while requiring repeated LLM calls and verifier-feedback rounds. This confirms that iterative repair is highly effective, but its gain comes with substantial LLM-iteration cost.

By contrast, symbolic-execution guidance provides a more cost-effective source of improvement. Even without refinement, adding symbolic-execution information raises \emph{Valid} over Pure LLM on every model: from \(13.0\%\) to \(31.5\%\) on gpt-4o, \(19.0\%\) to \(26.8\%\) on gpt-5-mini, \(26.5\%\) to \(33.0\%\) on gpt-5.4-mini, and \(13.8\%\) to \(36.0\%\) on gpt-5, whose unusually low \(13.8\%\) baseline is syntactic rather than semantic---it favors advanced ACSL constructs (behaviors, axiomatics) but often violates their surface syntax, triggering fatal whole-file parse errors that the symbolic-execution context helps avoid. Unlike refinement, this gain does not rely on repeated LLM sampling; it comes from exposing the local symbolic context before synthesis. Moreover, for loop-free programs fully supported by the flattened memory model, symbolic execution can directly compute the path-sensitive postcondition, so the core functional summary can be produced without an LLM call.

\rev{The refinement-budget ablation on \texttt{gpt-5.4-mini}
(\appref{app:refine-ablation}, Fig.~\supref{fig:refine-ablation}) shows that
the gain from refinement saturates after \(k{\approx}5\) iterations, that the
final Houdini pass contributes a further \(+14.1\)\,pp, and that the differential
refinement prompts with best-candidate rollback together account for roughly
23\,pp of additional validity over a plain error-plus-code repair prompt.}

\begin{finding}
\textbf{Finding (RQ5).} On \textbf{\toolname-400} the full pipeline reaches $68$--$93\%$ \emph{Valid}; refinement is the dominant but most expensive source of gain, whereas symbolic-execution guidance is a cheaper one that needs no repeated LLM sampling.
\end{finding}

\subsection{RQ6: How strong are \toolname's specifications versus the baselines?}
\label{subsec:rq6-strength}

\rev{Validity certifies that a specification is discharged by its verifier, but it does not rank how informative it is. We therefore compare accepted input domains (preconditions), normal-return guarantees (postconditions), and predicates at corresponding loop points (loop invariants) independently.}

\smallskip
\noindent\rev{\textbf{LLM judge and expert review.}}
\rev{Separate judge calls compare preconditions by accepted-domain inclusion, postconditions by bidirectional implication over their common input domain, and invariants at corresponding loop points. In Fig.~\ref{fig:fc_vs_judge_strength}, a direction means broader input acceptance or a stronger predicate, respectively; there is no aggregate ranking. Missing clauses mean \texttt{true}; unestablished implications count as Incomp. Expert review corrected 32 verdicts. \appref{app:strength} reports these corrections, prompts, reviewer details, and an A/B-order audit.}

\rev{Fig.~\ref{fig:fc_vs_judge_strength} shows stronger \toolname\ postconditions and loop invariants. Whereas \autospec\ accumulates sampled, verifier-accepted clauses, including redundant or tautological ones, \toolname\ supplies local symbolic states to guide path-sensitive and inductive predicates.}

\rev{\autospec\ more often declares broader input domains, commonly through a missing or \texttt{true} precondition; \toolname\ explicitly records runtime and memory safety requirements. This measures declared permissiveness, whose value depends on whether the additional inputs have meaningful valid executions.}

\rev{The \specgen\ comparison appears in \appref{app:strength} as cross-language context; the verifier stacks differ.}

\smallskip
\noindent\rev{\textbf{Frama-C/WP cross-check.}}
\rev{Separate Frama-C/WP harnesses cross-check 698 pre/postcondition pairs
and 540 invariant pairs using Frama-C 27.1, Alt-Ergo and Z3, with five
seconds per goal. In \toolname/baseline/equal/incomparable order, the counts
for preconditions, postconditions, and invariants are 90/281/302/25,
297/7/310/84, and 56/47/22/415, respectively.}

\rev{WP assigns a direction only when it proves every corresponding
obligation. An unproved implication enters Incomp whether it is false or true
but beyond the prover; unsupported harnesses and unresolved loop
correspondence also enter Incomp. Timeouts and automation limits therefore
leave many invariants unordered. The judge followed by expert review can
inspect semantic relations in these unresolved pairs, while WP provides an
independent formal check. Fig.~\ref{fig:fc_vs_judge_strength} confirms the
precondition and postcondition trends, but WP orders fewer invariants.}

\begin{figure}[t]
\centering
\rev{\footnotesize \tikz\filldraw[fill=figcontractfill,draw=figcontractstroke] (0,0) rectangle (.16,.16); \toolname\ side\quad
\tikz\filldraw[fill=figgrayfill,draw=figgrayframe] (0,0) rectangle (.16,.16); \autospec\ side\\[-1pt]
\tikz\filldraw[fill=figloopfill,draw=figloopstroke] (0,0) rectangle (.16,.16); Equal\quad
\tikz\filldraw[fill=figgray!10,draw=figgrayframe] (0,0) rectangle (.16,.16); Incomparable}

\newcommand{\strengthbar}[6]{%
  \node[anchor=east,text=\revcolor,font=\footnotesize] at (-2,#1) {#2};
  \path[draw=figcontractstroke,line width=.4pt,fill=figcontractfill] (0,#1-.24) rectangle (#3,#1+.24);
  \path[draw=figgrayframe,line width=.4pt,fill=figgrayfill] (#3,#1-.24) rectangle (#4,#1+.24);
  \path[draw=figloopstroke,line width=.4pt,fill=figloopfill] (#4,#1-.24) rectangle (#5,#1+.24);
  \path[draw=figgrayframe,line width=.4pt,fill=figgray!10] (#5,#1-.24) rectangle (100,#1+.24);
  \node[anchor=west,text=\revcolor,font=\footnotesize] at (102,#1) {#6};}
\begin{tikzpicture}[x=.45mm,y=.46cm,every node/.style={font=\footnotesize}]
  \path[use as bounding box] (-45,.05) rectangle (147,8.45);
  \foreach \x in {0,25,50,75,100} {
    \draw[figgrayframe!70,densely dotted] (\x,.65) -- (\x,7.75);
    \node[anchor=north,text=\revcolor,font=\footnotesize] at (\x,.55) {\x\%};}

  \node[anchor=west,text=\revcolor,font=\footnotesize\bfseries] at (0,8.08) {(a) Preconditions ($N=698$)};
  \strengthbar{7.25}{Judge + expert}{13.75}{55.30}{99.86}{96/290/311/1}
  \strengthbar{6.30}{WP}{12.89}{53.15}{96.42}{90/281/302/25}

  \node[anchor=west,text=\revcolor,font=\footnotesize\bfseries] at (0,5.38) {(b) Postconditions ($N=698$)};
  \strengthbar{4.55}{Judge + expert}{49.00}{51.58}{97.42}{342/18/320/18}
  \strengthbar{3.60}{WP}{42.55}{43.55}{87.97}{297/7/310/84}

  \node[anchor=west,text=\revcolor,font=\footnotesize\bfseries] at (0,2.68) {(c) Loop invariants ($N=540$)};
  \strengthbar{1.85}{Judge + expert}{22.22}{31.85}{32.78}{120/52/5/363}
  \strengthbar{.90}{WP}{10.37}{19.07}{23.15}{56/47/22/415}
\end{tikzpicture}
\caption{\rev{Dimension-wise relations for \autospec\ versus \toolname\ over identical data sets. A directional segment means a more permissive accepted domain for preconditions and a logically stronger predicate for postconditions or invariants. ``Judge + expert'' is the corrected judge result; WP checks the same implications. Right-hand counts follow the legend order.}}
\label{fig:fc_vs_judge_strength}
\end{figure}
\let\strengthbar\relax

\begin{finding}
\rev{\textbf{Finding (RQ6).} Symbolic context changes where specification strength appears: \toolname\ yields more informative postconditions and stronger loop invariants than \autospec\ in the shared C/ACSL setting, while \autospec's broader declared input domains mainly reflect safety assumptions that are not written explicitly. The cross-language comparison against \specgen\ (Java/JML) is reported in \appref{app:strength}; no single total ordering applies across all dimensions and languages.}
\end{finding}


\subsection{RQ7: What does \toolname\ cost in time, tokens, API calls, and memory?}
\label{subsec:rq7-cost}

Costs are measured on the 400 goal-independent programs, with all methods
under Pass@1. \rev{Fig.~\ref{fig:rq7-cost} compares every attempt, including
failures, on \texttt{gpt-5.4-mini}; the original \autospec\ time, call count,
and validity labels are retained.} Peak RSS stays below \(\sim 500\)~MB.

\begin{figure*}[t]
\centering
\begin{tikzpicture}[x=1cm,y=0.50cm,font=\footnotesize,every node/.style={text=figgraytext}]
  \def\coststack#1#2#3#4#5#6#7{%
    \path[fill=#5] (#1,{#2-0.18}) rectangle (#4,{#2+0.18});
    \path[pattern=north east lines,pattern color=figgraytext!65]
      (#3,{#2-0.18}) rectangle (#4,{#2+0.18});
    \draw[draw=#6,line width=0.45pt] (#1,{#2-0.18}) rectangle (#4,{#2+0.18});
    \draw[draw=#6,line width=0.35pt] (#3,{#2-0.18}) -- (#3,{#2+0.18});
    \node[anchor=west] at ({#4+0.05},#2) {\rev{#7}};
  }

  \path[fill=figcontractfill] (4.55,5.88) rectangle (4.85,6.18);
  \draw[draw=figcontractstroke] (4.55,5.88) rectangle (4.85,6.18);
  \node[anchor=west] at (4.91,6.03) {\rev{Successful runs}};
  \path[fill=figcontractfill] (7.30,5.88) rectangle (7.60,6.18);
  \path[pattern=north east lines,pattern color=figgraytext!65]
    (7.30,5.88) rectangle (7.60,6.18);
  \draw[draw=figcontractstroke] (7.30,5.88) rectangle (7.60,6.18);
  \node[anchor=west] at (7.66,6.03) {\rev{Failed runs}};

  \node[anchor=east] at (-0.28,4) {\autospec};
  \node[anchor=east] at (-0.28,3) {\specgen};
  \node[anchor=east] at (-0.28,2) {Pure LLM};
  \node[anchor=east] at (-0.28,1) {\toolname\ w/o refine};
  \node[anchor=east] at (-0.28,0) {\textbf{\toolname}};

  \node[anchor=south] at (1.9,4.55) {Time (s)};
  \draw[figaxis] (0,-0.45) -- (3.8,-0.45);
  \node[anchor=north] at (0,-0.55) {0};
  \node[anchor=north] at (3.8,-0.55) {\rev{1200}};
  \node[anchor=south] at (6.8,4.55) {Tokens};
  \draw[figaxis] (4.9,-0.45) -- (8.7,-0.45);
  \node[anchor=north] at (4.9,-0.55) {0};
  \node[anchor=north] at (8.7,-0.55) {120k};
  \node[anchor=south] at (11.7,4.55) {API calls};
  \draw[figaxis] (9.8,-0.45) -- (13.6,-0.45);
  \node[anchor=north] at (9.8,-0.55) {0};
  \node[anchor=north] at (13.6,-0.55) {15};

  % Solid segment = successful cost / all attempts; hatched = failed cost / all attempts.
  \coststack{0.000}{4}{1.307}{2.952}{figgrayfill}{figgrayframe}{932.1}
  \coststack{0.000}{3}{0.445}{1.189}{figgrayfill}{figgrayframe}{375.5}
  \coststack{0.000}{2}{0.008}{0.039}{figgrayfill}{figgrayframe}{12.3}
  \coststack{0.000}{1}{0.051}{0.129}{figloopfill}{figloopstroke}{40.8}
  \coststack{0.000}{0}{0.856}{1.008}{figcontractfill}{figcontractstroke}{318.2}

  \coststack{4.900}{4}{5.789}{6.297}{figgrayfill}{figgrayframe}{44.1k}
  \coststack{4.900}{3}{4.966}{5.039}{figgrayfill}{figgrayframe}{4.4k}
  \coststack{4.900}{2}{4.913}{4.952}{figgrayfill}{figgrayframe}{1.6k}
  \coststack{4.900}{1}{4.955}{5.129}{figloopfill}{figloopstroke}{7.2k}
  \coststack{4.900}{0}{7.166}{7.553}{figcontractfill}{figcontractstroke}{83.8k}

  \coststack{9.800}{4}{11.859}{13.072}{figgrayfill}{figgrayframe}{12.9}
  \coststack{9.800}{3}{10.022}{10.269}{figgrayfill}{figgrayframe}{1.9}
  \coststack{9.800}{2}{9.867}{10.053}{figgrayfill}{figgrayframe}{1.0}
  \coststack{9.800}{1}{9.927}{10.304}{figloopfill}{figloopstroke}{2.0}
  \coststack{9.800}{0}{11.625}{11.896}{figcontractfill}{figcontractstroke}{8.3}
\end{tikzpicture}

\caption{\rev{Mean cost per attempted program on all 400 \texttt{gpt-5.4-mini} cases. Each bar stacks valid-run (solid) and failed-run (hatched) contributions, each divided by 400.}}
\label{fig:rq7-cost}
\end{figure*}

\rev{Each solid (hatched) bar segment sums valid-run (failed-run) costs over
all 400 attempts. Failed runs contribute 48.1 of \toolname's 318.2 seconds
per attempt, versus 519.4/932.1 for \autospec\ and 235.0/375.5 for
\specgen. Time per obtained valid specification is 338.6, 1,462.1, and
758.6 seconds, respectively. \toolname's 83.8k tokens per attempt include
repeated refinement context without cache discounts.}

\begin{finding}
\textbf{Finding (RQ7).} \rev{\toolname's added framework complexity mainly appears as higher token use; its wall-clock time and API-call count are much closer to the baselines than the raw token totals alone suggest.}
\end{finding}


\subsection{RQ8: What are \toolname's dominant failure modes?}
\label{subsec:rq8-failure}

Fig.~\ref{fig:failure_root_causes} analyzes an independent
\texttt{gpt-5.4-mini} rerun with per-case logs. Capping refinement at three
iterations instead of ten exposes more failures: 321 valid cases versus 373
in the main run. The panels show terminal outcomes and the dominant program
feature in each of the 79 non-valid cases (42 \emph{No valid}, 37 \emph{No
syntax}); for overlapping features, we choose the first explaining the
blocking obligation. Table~\supref{tab:gpt54mini_failure_events}
(\appref{sec:limitations}) separately classifies logged syntax and prover
events by outcome and obligation type.

Loop invariant obligations account for \(2{,}590/4{,}370=59.3\%\) of
timeouts, especially establishment and preservation over arrays and heap
arrays. Postcondition and \texttt{assigns} timeouts also occur.

\rev{Log outcomes \emph{Timeout}, \emph{Unknown}, and \emph{Failed} denote
timeouts, prover uncertainty, and backend or resource failures, not semantic
counterexamples. They cannot separate incorrect specifications from correct
but unproved obligations. Manual review more often found plausible properties
limited by complex loop, array, or heap proofs, without ruling out synthesis
errors.} See \rev{\appref{sec:limitations}} for method and verifier limits.

\begin{figure}[t]
\centering
\begin{minipage}[t]{\columnwidth}
\centering
\begin{tikzpicture}[font=\footnotesize,every node/.style={text=figgraytext}]
  \path[use as bounding box] (-1.20,-1.10) rectangle (6.25,1.10);
  \def\r{1.00}
  \filldraw[draw=figgrayframe,line width=0.3pt,fill=figgrayfill]
    (0,0) -- (0:\r) arc[start angle=0,end angle=288.9,radius=\r] -- cycle;
  \filldraw[draw=figcontractstroke,line width=0.3pt,fill=figcontractfill]
    (0,0) -- (288.9:\r) arc[start angle=288.9,end angle=326.7,radius=\r] -- cycle;
  \filldraw[draw=figcallstroke,line width=0.3pt,fill=figcallfill]
    (0,0) -- (326.7:\r) arc[start angle=326.7,end angle=360,radius=\r] -- cycle;
  \draw[figgrayframe,line width=0.3pt] (0,0) circle (\r);
  \filldraw[fill=figgrayfill,draw=figgrayframe,line width=0.25pt] (1.55,0.54) rectangle +(0.16,0.16);
  \node[anchor=west] at (1.82,0.62) {Valid: 321 (80.3\%)};
  \filldraw[fill=figcontractfill,draw=figcontractstroke,line width=0.25pt] (1.55,0.22) rectangle +(0.16,0.16);
  \node[anchor=west] at (1.82,0.30) {No valid: 42 (10.5\%)};
  \filldraw[fill=figcallfill,draw=figcallstroke,line width=0.25pt] (1.55,-0.10) rectangle +(0.16,0.16);
  \node[anchor=west] at (1.82,-0.02) {No syntax: 37 (9.2\%)};
\end{tikzpicture}
\par\smallskip{\footnotesize\normalfont (a) Terminal outcomes.}
\end{minipage}
\par\smallskip
\begin{minipage}[t]{\columnwidth}
\centering
\begin{tikzpicture}[font=\footnotesize,every node/.style={text=figgraytext}]
  \path[use as bounding box] (-1.20,-1.15) rectangle (6.25,1.15);
  \def\r{1.00}
  \filldraw[draw=figgrayframe,line width=0.28pt,fill=figgrayfill]
    (0,0) -- (0:\r) arc[start angle=0,end angle=141.27,radius=\r] -- cycle;
  \filldraw[draw=figloopstroke,line width=0.28pt,fill=figloopfill]
    (0,0) -- (141.27:\r) arc[start angle=141.27,end angle=250.63,radius=\r] -- cycle;
  \filldraw[draw=figgrayframe,line width=0.28pt,fill=figgray!10]
    (0,0) -- (250.63:\r) arc[start angle=250.63,end angle=314.43,radius=\r] -- cycle;
  \filldraw[draw=figcallstroke,line width=0.28pt,fill=figcallfill]
    (0,0) -- (314.43:\r) arc[start angle=314.43,end angle=341.77,radius=\r] -- cycle;
  \filldraw[draw=figloopstroke,line width=0.28pt,fill=figloop!20]
    (0,0) -- (341.77:\r) arc[start angle=341.77,end angle=350.89,radius=\r] -- cycle;
  \filldraw[draw=figgrayframe,line width=0.28pt,fill=figgray!6]
    (0,0) -- (350.89:\r) arc[start angle=350.89,end angle=355.44,radius=\r] -- cycle;
  \filldraw[draw=figgrayframe,line width=0.28pt,fill=white]
    (0,0) -- (355.44:\r) arc[start angle=355.44,end angle=360,radius=\r] -- cycle;
  \draw[figgrayframe,line width=0.3pt] (0,0) circle (\r);
  \filldraw[fill=figgrayfill,draw=figgrayframe,line width=0.2pt] (1.45,0.895) rectangle +(0.13,0.13);
  \node[anchor=west] at (1.68,0.96) {Heap dyn. arrays: 31 (39.2\%)};
  \filldraw[fill=figloopfill,draw=figloopstroke,line width=0.2pt] (1.45,0.575) rectangle +(0.13,0.13);
  \node[anchor=west] at (1.68,0.64) {Array pointer loops: 24 (30.4\%)};
  \filldraw[fill=figgray!10,draw=figgrayframe,line width=0.2pt] (1.45,0.255) rectangle +(0.13,0.13);
  \node[anchor=west] at (1.68,0.32) {String char logic: 14 (17.7\%)};
  \filldraw[fill=figcallfill,draw=figcallstroke,line width=0.2pt] (1.45,-0.065) rectangle +(0.13,0.13);
  \node[anchor=west] at (1.68,0.00) {Recursive callee assert: 6 (7.6\%)};
  \filldraw[fill=figloop!20,draw=figloopstroke,line width=0.2pt] (1.45,-0.385) rectangle +(0.13,0.13);
  \node[anchor=west] at (1.68,-0.32) {Scalar loops: 2 (2.5\%)};
  \filldraw[fill=figgray!6,draw=figgrayframe,line width=0.2pt] (1.45,-0.705) rectangle +(0.13,0.13);
  \node[anchor=west] at (1.68,-0.64) {Floating-point casts: 1 (1.3\%)};
  \filldraw[fill=white,draw=figgrayframe,line width=0.2pt] (1.45,-1.025) rectangle +(0.13,0.13);
  \node[anchor=west] at (1.68,-0.96) {Other: 1 (1.3\%)};
\end{tikzpicture}
\par\smallskip{\footnotesize\normalfont (b) Root causes for non-valid cases.}
\end{minipage}
\caption{Failure outcome and root-cause summary for the \texttt{gpt-5.4-mini} \toolname-400 rerun.}
\label{fig:failure_root_causes}
\end{figure}

\begin{finding}
\textbf{Finding (RQ8).} \rev{Residual failures manifest predominantly as SMT timeouts on complex proof obligations. The logs cannot conclusively separate specification-synthesis errors from proof obligations beyond the current backend, although manual inspection more often points to verifier limitations.}
\end{finding}

